% Generated draft: requires mathematical and source review.
\documentclass[11pt]{article}
\usepackage[margin=25mm]{geometry}
\usepackage{fontspec}
\setmainfont{Latin Modern Roman}
\usepackage{amsmath,amssymb,mathtools,unicode-math}
\setmathfont{Latin Modern Math}
\usepackage{enumitem,needspace,xurl,hyperref}
\setlength{\parindent}{0pt}
\setlength{\parskip}{5pt}
\setlength{\emergencystretch}{3em}
\newcommand{\N}{\mathbb{N}}
\newcommand{\Z}{\mathbb{Z}}
\newcommand{\Q}{\mathbb{Q}}
\newcommand{\R}{\mathbb{R}}
\newcommand{\C}{\mathbb{C}}
\newcommand{\F}{\mathbb{F}}
\newcommand{\A}{\mathbb{A}}
\newcommand{\PP}{\mathbb{P}}
\newcommand{\E}{\mathbb{E}}
\newcommand{\eps}{\varepsilon}
\newcommand{\dd}{\,\mathrm{d}}
\begin{document}
\section*{UnsolvedMath ID: 20002487}
\textbf{Source problem code:} catalogue entry
\subsection{Definitions and mathematical statement}
Craig Evans asks whether Mather sets play a role in semiclassical eigenfunctions for $H(x,p)=\tfrac12|p|^2+V(x)$ with periodic potential, and conversely whether Mather sets can be quantized by quasimodes with good error bounds; he also asks whether a Diophantine condition would be useful. The aid identifies this as original Problem 31, whose leading digit was lost in the extraction. The source first gives the operator $-h^2\Delta/2+V$ but prints the opposite sign in its eigenvalue display. With $\Delta=\sum_j\partial_{x_j}^2$, the Hamiltonian-consistent choice used by the aid is
\[
P_h=-\tfrac{h^2}{2}\Delta+V,\qquad h\downarrow0.
\]
The sign discrepancy is retained as a source qualification rather than silently resolved for both printed formulas. Use a smooth real periodic $V$ for the following supplied claims; the original does not specify its regularity.

The mechanical Lagrangian is $L(x,v)=\tfrac12|v|^2-V(x)$. At zero cohomology, put $V_{\max}=\max V$ and $K=\{x:V(x)=V_{\max}\}$. The aid asserts that all minimizing invariant probability measures are exactly
\[
\mu(\dd x)\delta_0(\dd v),\qquad \mu\text{ any probability measure supported on }K,
\]
with momentum $p=v$. In particular the zero-class Mather set is $K\times\{0\}$. It reports that every such measure is realized by normalized smooth periodic quasimodes along a sequence $h_j\downarrow0$ with
\[
\|(P_{h_j}-V_{\max})u_{h_j}\|_{L^2}\leq C h_j,\qquad
\|u_{h_j}\|_{L^2}=1.
\]
Precisely, their semiclassical Wigner distributions satisfy
\[
\langle\operatorname{Op}_{h_j}(a)u_{h_j},u_{h_j}\rangle
\longrightarrow\int_K a(x,0)\,\dd\mu(x)
\]
for standard smooth compact-momentum test symbols. The proposed construction is a diagonal sequence of superpositions of separated packets of width $\sqrt h$, with a residual constant uniform in the growing packet count. It includes singular and nonatomic measures on degenerate maximum sets. This is an $O(h)$ quasimode statement, not an exact-eigenfunction realization or an arbitrary-accuracy claim.

The second supplied construction assumes a smooth global Hamilton--Jacobi graph represented locally by $\eta=\dd S$, where $\eta$ is a closed real one-form on the torus and $H(x,\eta(x))=E$. Assume a positive smooth normalized invariant density $\rho$ with $\operatorname{div}(\rho\eta)=0$, and set $a=\sqrt\rho$. Whenever the global phase satisfies the Bohr--Sommerfeld condition
\[
\frac1h\int_\gamma\eta\in2\pi\Z\quad\text{for every closed torus loop }\gamma,
\]
the phase $e^{iS/h}$ is globally periodic and
\[
u_h=a e^{iS/h},\qquad
(P_h-E)u_h=-\tfrac{h^2}{2}e^{iS/h}\Delta a.
\]
Thus the residual is $O(h^2)$ for this fixed smooth graph and amplitude. The condition applies to admissible values of $h$ and does not assert that an admissible sequence exists for every cohomology class. The packet result needs no Diophantine condition; the graph result uses the stated holonomy condition, which is not a general answer to the source's broader arithmetic question. Neither construction treats every nonzero-cohomology or nonsmooth Mather set.
\subsection{Short English statement}
Evans asks how Mather sets could be realized semiclassically. The aid realizes every zero-class minimizing measure over potential maxima by periodic $O(h)$ quasimodes, and gives $O(h^2)$ WKB states on smooth invariant graphs when the global phase is quantized. These partial constructions have distinct regularity, cohomology, and arithmetic restrictions.
\subsection{Sources}
\begin{itemize}
\item Local UnsolvedMath registry, permanent ID $20002487$, \texttt{catalogue entry}; \texttt{aim-pdes-notes.json}, \texttt{notes[98]}, extracted source item 1, identified by supplied verification as Problem 31.
\item American Institute of Mathematics, \emph{New connections between dynamical systems and PDE's}, Craig Evans contribution, \url{https://aimath.org/WWN/dynpde/dynpde.pdf}. The isolated $3$ is a page artifact. No authoritative separate bibliographic reference for the realization or WKB claims is supplied.
\item Research attempt 1: partial result, problem partially solved at run, rigor label ``proved,'' novelty confidence low. Registry classification partial progress, status partially solved, statement status exact. The source-verification field reports the numbering correction but no full recovered sign-corrected display.
\end{itemize}
\end{document}
